\newpage
\begin{abstract}
Building on the general connection between culture consumption and sociability which has been the focus of recent research in the sociology of culture, I propose that persons who have a propensity to consume a wide variety of cultural goods should also be observed to have successfully mobilized network ties for instrumental purposes, especially ``weak ties'' in contrast to ``strong ties.'' The results show that net of other individual predictors, as the cultural variety of the person's repertoire of tastes increases, there is 1) a higher likelihood of having activated a social tie to obtain information about jobs, and 2) a higher likelihood that the tie activation is through a weak, rather than strong ties or non-relational mechanisms.
\end{abstract}


\begin{center}
Word Count: 4,092
\end{center}